%!TeX root=MemoriaTFG.tex

\chapter{Introducción}
\label{cap:intro}

\section{Contexto clínico}
\label{sec:contexto}

La dermatología clínica atiende un espectro amplio de entidades con
distribución concentrada en un número reducido de motivos de
consulta. El estudio observacional DIADERM, sobre $10\,999$
diagnósticos en $8\,953$ pacientes en territorio español,
ordena los diez diagnósticos más prevalentes como queratosis
actínica, carcinoma basocelular, nevus melanocítico, queratosis
seborreica, otras neoplasias benignas, psoriasis, acné, verrugas
víricas, otros trastornos de la pigmentación y otras
dermatitis~\cite{diaderm2018}. El trabajo local de Taberner et~al.\
sobre el servicio dermatológico del Hospital Universitario Son
Llàtzer (HUSLL) documenta $213$ diagnósticos diferentes en un
único año asistencial~\cite{taberner2010_motivos}. Ambos coinciden
en un peso elevado de dermatosis benignas, inflamatorias e
infecciosas habituales.

La sensibilidad temporal del proceso diagnóstico resulta crítica en
el melanoma cutáneo. La supervivencia específica a cinco años
desciende desde aproximadamente $99\%$ en estadio~I hasta
$25\%$ en estadio~IV según la octava edición de la American
Joint Committee on Cancer~\cite{ajcc8}. El acceso al especialista
se realiza mediante derivación desde atención primaria. La
capacidad asistencial del circuito presencial es limitada frente a
una demanda creciente, lo que justifica el interés por sistemas de
apoyo al diagnóstico basados en imagen.

El sustrato visual del diagnóstico dermatológico se articula en tres
modalidades: fotografía clínica (presentación general y contexto
anatómico), dermatoscopia (patrones pigmentarios y estructurales no
visibles a simple vista) e histopatología (confirmación
tisular). Sobre esta tríada operan los sistemas de apoyo basados en
aprendizaje automático.

\section{Análisis automatizado: paradigma clásico y limitaciones}
\label{sec:paradigma-clasico}

El primer hito establecido del aprendizaje profundo en dermatología
es el clasificador convolucional de Esteva et~al.~\cite{esteva2017},
con rendimiento equivalente al de dermatólogos certificados sobre
un conjunto curado de lesiones. La aproximación clásica entrena
arquitecturas convolucionales desde inicialización aleatoria o desde
pesos preentrenados en ImageNet, ajustadas posteriormente sobre
datasets dermatológicos etiquetados~\cite{tschandl2018,ham10000}. La
literatura posterior incorpora aumento de datos, balanceo de clases,
ensembles y arquitecturas específicas para tareas concretas como
segmentación binaria~\cite{isic2018} o clasificación binaria
melanoma/no-melanoma~\cite{kawahara2019}.

Tres limitaciones definen este paradigma. La primera es el coste de
anotación: el rendimiento escala con el volumen de datos etiquetados
y la anotación dermatológica requiere expertos clínicos con
disponibilidad limitada. La segunda es la fragmentación entre
conjuntos públicos: los datasets disponibles utilizan taxonomías
heterogéneas, distintos protocolos de adquisición y diferentes
esquemas de balance de clases, lo que limita comparaciones directas
entre estudios y transferencia entre dominios. La tercera es la
pérdida de rendimiento fuera del dominio de entrenamiento,
documentada de forma reiterada en pacientes de fototipo cutáneo
elevado~\cite{fitzpatrick17k,ddi}. Estas tres limitaciones motivan
la transición hacia el paradigma fundacional.

\section{Modelos fundacionales en imagen médica}
\label{sec:fundacionales-medica}

Los modelos fundacionales son arquitecturas preentrenadas sobre
grandes datasets sin etiquetar mediante objetivos auto-supervisados o
multimodales. Sus representaciones se adaptan posteriormente a
tareas específicas con cantidades reducidas de datos etiquetados.
Las arquitecturas predominantes son el Vision Transformer (ViT) y
sus variantes~\cite{vit2021}. Los objetivos de preentrenamiento más
extendidos son la reconstrucción enmascarada, la destilación
auto-supervisada y la alineación contrastiva imagen-texto al estilo
CLIP~\cite{clip2021,siglip2023}.

En imagen médica, DINOv2~\cite{dinov2} representa la referencia
abierta de encoder auto-supervisado de propósito general, entrenado
sobre $\sim 142$ millones de imágenes naturales mediante
destilación entre vistas. BiomedCLIP~\cite{biomedclip} extiende el
paradigma contrastivo al dominio biomédico sobre PMC-15M
($\sim 15$ millones de pares figura-leyenda). El paradigma de
segmentación promptable se establece con SAM~\cite{sam2023} y su
evolución SAM2.1~\cite{sam2}; su transferencia a dominios médicos
requiere adaptación, habitualmente vía LoRA~\cite{lora}. La
arquitectura BLIP-2~\cite{blip2_2023} bridge un encoder visual
congelado con un modelo de lenguaje mediante un módulo
intermedio (Q-Former). La evidencia disponible muestra que el
preentrenamiento sobre datos del dominio clínico de interés
incrementa de forma consistente el rendimiento respecto a encoders
generalistas adaptados.

\section{Modelos fundacionales para dermatología}
\label{sec:fundacionales-derma}

Tres referencias relevantes liberadas por el grupo de la Universidad
Monash definen el estado del arte específico durante el periodo
2024--2026.

\paragraph{PanDerm~\cite{panderm2025}.}
Encoder visual auto-supervisado entrenado mediante reconstrucción
enmascarada sobre $\sim 2{,}1$ millones de imágenes
dermatológicas en cuatro modalidades (fotografía corporal total,
dermatoscopia, fotografía clínica, histopatología). Libera pesos en
dos tamaños: \emph{Base} ($\sim 86$\,M parámetros, embedding
$768$) y \emph{Large} ($\sim 307$\,M parámetros, embedding
$1\,024$). Sus autores reportan una mejora media del
$+10{,}2\%$ sobre el estado del arte previo en $28$ tareas
dermatológicas. La disponibilidad simultánea de pesos, código y
benchmark hace de PanDerm el primer modelo fundacional dermatológico
verificable de forma independiente.

\paragraph{Derm1M y la familia DermLIP~\cite{derm1m}.}
Derm1M es un dataset dermatológico vision-lenguaje de
$\sim 1{,}03$ millones de pares imagen-texto, construido a
partir de fuentes clínicas y de literatura especializada. DermLIP es
el modelo contrastivo entrenado sobre Derm1M al estilo
CLIP~\cite{clip2021}, distribuido en variantes con encoder visual
ViT-B/16 y encoder textual basado en PubMedBERT o GPT-2/CLIP. La
alineación habilita clasificación zero-shot guiada por
\emph{prompts} y recuperación bidireccional visión-texto sin
entrenamiento específico por tarea.

\paragraph{DermFM-Zero~\cite{dermfmzero2026}.}
Integra encoder visual, rama lingüística y mecanismo de
interpretabilidad mediante Sparse Autoencoders sobre las
representaciones internas~\cite{templeton2024}. Reporta rendimientos
superiores a PanDerm y DermLIP en los benchmarks que define. Sus
pesos no son públicos en la fecha de cierre del trabajo, lo que
restringe su tratamiento a la discusión conceptual.

A estos tres modelos específicos se añaden, en el espacio de
comparación, encoders generalistas (DINOv2, BiomedCLIP,
SigLIP-Large), arquitecturas convolucionales clásicas (ConvNeXt-L,
EfficientNetV2-L), el modelo de segmentación SAM2.1 con adaptación
LoRA y modelos de lenguaje multimodales (GPT-4o, Gemini~2.5~Pro,
MedGemma~4B Vision, MedGemma~27B Text)~\cite{amie2024}, junto con
BLIP-2 como referencia metodológica.

\section{Brechas operativas y motivación}
\label{sec:brechas}

La revisión de la literatura identifica cinco brechas operativas
sobre las que el presente estudio aporta evidencia empírica.

\paragraph{Benchmarks no homogéneos entre familias.}
Los benchmarks reportados por los autores originales de PanDerm,
DermLIP y DermFM-Zero se construyen, en buena medida, sobre
subconjuntos derivados de los datos de preentrenamiento, y aplican
protocolos no equivalentes entre familias. La comparación entre
encoders específicos de dominio, modelos contrastivos
vision-lenguaje, encoders generalistas y modelos de lenguaje
multimodales sobre un mismo conjunto de datasets externos y un
mismo protocolo no está documentada en la literatura disponible.

\paragraph{Heterogeneidad taxonómica entre datasets.}
Los doce datasets externos públicos identificados como accesibles
para el trabajo utilizan taxonomías nativas incompatibles:
binarias melanoma/no-melanoma en Derm7pt, HIBA y MSKCC;
multiclase con $C\in[6,23]$ en HAM10000, BCN20000, PAD-UFES-20,
Dermnet y WSI~patches; $114$ patologías en
Fitzpatrick17k~\cite{fitzpatrick17k}; $48$ conceptos clínicos en
SkinCon~\cite{skincon}. La comparación directa entre modelos sobre
estos datasets exige un esquema explícito de armonización
taxonómica que la literatura no proporciona.

\paragraph{Fuga de información entre conjuntos de evaluación y
preentrenamiento.}
La construcción de Derm1M a partir de fuentes web abiertas
introduce un riesgo no cuantificado de solapamiento con los
conjuntos públicos habituales de evaluación. La cuantificación
sistemática mediante hash MD5 exhaustivo no se documenta en la
literatura previa al inicio del trabajo y constituye condición
necesaria para la interpretación defendible de cualquier resultado
zero-shot sobre estos datasets.

\paragraph{Datasets dermatológicos en español.}
Los datasets disponibles para preentrenamiento o adaptación de
modelos vision-lenguaje son mayoritariamente angloparlantes
(Derm1M, MONET, MM-Skin). La disponibilidad de material clínico en
español con calidad docente verificada (imagen, metadatos clínicos
y texto narrativo asociado) es marginal.

\paragraph{Reproducibilidad sobre hardware no estándar.}
La práctica totalidad de la literatura experimental se ejecuta
sobre GPU NVIDIA A100 o H100 en plataformas x86\_64. La aparición
de plataformas alternativas basadas en arquitectura aarch64 (ARM de
64 bits), como la familia NVIDIA DGX Spark con chip Grace Hopper,
introduce una variable no evaluada sobre la viabilidad de
reproducción y el comportamiento numérico de los pipelines en BF16.

\section{Marco operativo previo del autor}
\label{sec:trabajo-previo}

El presente trabajo se desarrolla en el marco de dos colaboraciones
preexistentes con la Dra.~R.~Taberner, dermatóloga del HUSLL. La
primera es la modernización tecnológica del archivo Dermapixel,
blog mantenido por la Dra.~Taberner desde 2010 que reúne más de
$700$ casos clínicos comentados con imagen clínica y
dermatoscópica, historia abreviada, diagnóstico y razonamiento
diferencial; el autor trabaja en su migración a una pila
tecnológica propia. La segunda es un proyecto institucional del
Servei de Salut de les Illes Balears (IB-Salut) sobre
teledermatología hospitalaria, en el que se desarrolla una
herramienta de apoyo a la primera valoración dermatológica en el
circuito de derivación desde atención primaria. Ambas líneas
orientan la decisión de priorizar modelos con código y pesos
públicos por su viabilidad de integración.

\section{Pregunta de investigación}
\label{sec:pregunta}

El estudio se articula en torno a una pregunta única: \emph{¿cuál
es el rendimiento y la capacidad de generalización de los modelos
fundacionales para imagen dermatológica disponibles, en función del
volumen de datos etiquetados y del tipo de supervisión, y con qué
robustez frente a la heterogeneidad taxonómica entre datasets y la
disparidad por fototipo cutáneo?}

La respuesta empírica exige dos condiciones operativas: un
protocolo de evaluación homogéneo aplicable a familias de modelos
heterogéneas y un esquema explícito de armonización taxonómica
entre datasets. El diseño descrito en el
capítulo \emph{Materiales y métodos} (en preparación) satisface ambas.

\section{Objetivos del trabajo}
\label{sec:objetivos}

El trabajo se articula en torno a dos objetivos formales,
verificables sobre los resultados experimentales del
capítulo \emph{Resultados} (en preparación).

\paragraph{O1. Evaluación experimental sistemática de los modelos
fundacionales dermatológicos sobre conjuntos de datos externos.}
Sobre el hardware del proyecto (NVIDIA DGX Spark, chip Grace Hopper
GB10, arquitectura aarch64, $128$\,GB memoria unificada,
CUDA~13.0, BF16) se implementa el conjunto de modelos fundacionales
con código y pesos públicos a la fecha de inicio. Se evalúa su
comportamiento sobre doce datasets externos
públicos\footnote{HAM10000~\cite{ham10000},
BCN20000~\cite{bcn20000}, PAD-UFES-20~\cite{padufes},
DDI~\cite{ddi}, Dermnet~\cite{dermnet}, Derm7pt
(clínica y dermatoscópica)~\cite{kawahara2019}, HIBA, MSKCC, WSI
patches, Fitzpatrick17k~\cite{fitzpatrick17k},
SkinCon~\cite{skincon} e ISIC2018~\cite{isic2018}.} mediante cuatro
protocolos principales (sondeo lineal, ajuste fino supervisado,
segmentación binaria con LoRA~\cite{lora} sobre SAM2.1~\cite{sam2},
clasificación zero-shot multimodal) y tres protocolos
complementarios (comparación con modelos de lenguaje multimodales,
análisis de equidad por fototipo Fitzpatrick, auditoría de
solapamiento con Derm1M mediante hash MD5).

\paragraph{O2. Construcción del dataset DermapixelAI 1.0.}
Se extrae, depura y estructura el material del archivo Dermapixel
para conformar un dataset que vincula, para cada caso, imagen y
modalidad, metadatos clínicos y texto narrativo en español. El
dataset se publica bajo licencia CC-BY-NC-SA~4.0 con datasheet
formal~\cite{datasheets2018}, particiones documentadas y mecanismo
de cita recomendada. La explotación experimental sistemática sobre
el dataset se inicia hacia el cierre del periodo cubierto por este
documento y se desarrolla principalmente con posterioridad;
la apartado \emph{Líneas de investigación abiertas} (en preparación) del capítulo \emph{Conclusiones} (en preparación)
recoge el estado del dataset.

Los bloques derivados emergidos durante el desarrollo del proyecto
(armonización ontológica jerárquica L1/L2/L3, ensemble de seguridad
para detección de melanoma, análisis de equidad por fototipo,
interpretabilidad mediante Sparse Autoencoders~\cite{templeton2024},
recuperación visual densa basada en FAISS~\cite{faiss2019},
integración en el prototipo DermApIxel) no constituyen objetivos
formales y se documentan en los anexos correspondientes con su
estado, prerrequisitos y líneas futuras.

\section{Contribuciones}
\label{sec:contribuciones}

Las contribuciones se enumeran a continuación. Las cuatro primeras
sostienen la respuesta a la pregunta de investigación y se
desarrollan en el cuerpo principal; las cinco siguientes constituyen
material complementario y se documentan en los
Anexos F--H (en preparación).

\begin{enumerate}
\item Reproducción experimental del modelo PanDerm sobre
arquitectura aarch64 (GB10/Grace Hopper) con verificación numérica
respecto a los valores reportados en la publicación original.

\item Benchmark comparativo homogéneo de catorce modelos (PanDerm
Base y Large, DINOv2, ConvNeXt-Large, EfficientNetV2-Large, tres
variantes de DermLIP, BiomedCLIP, SigLIP-Large, SAM2.1+LoRA, GPT-4o,
Gemini~2.5~Pro, MedGemma~4B Vision y~27B Text, BLIP-2) bajo cuatro
protocolos sobre doce datasets externos.

\item Ontología jerárquica L1/L2/L3 para armonización taxonómica
entre datasets heterogéneos: cuatro categorías etiológicas,
$43$ subcategorías clínicas y $367$ diagnósticos
individuales, codificados en SNOMED~CT y CIE-10 y validados por
revisión experta.

\item Auditoría sistemática de solapamiento entre los doce
conjuntos de evaluación y el dataset Derm1M mediante hash MD5
exhaustivo, con identificación de tres datasets afectados
(Dermnet, HIBA, MSKCC) y cuantificación del impacto sobre las
cifras zero-shot.

\item Dataset DermapixelAI~1.0, primer dataset dermatológico
verificado en español con imagen, metadatos y texto narrativo
asociado ($1\,089$ imágenes vinculadas a $672$ casos clínicos,
mediana de $223$ palabras por caso), distribuido bajo licencia
CC-BY-NC-SA~4.0 con datasheet, particiones \emph{case-aware}
($891 / 157 / 41$) y mecanismo de cita
(Anexo \emph{Entrega del dataset} (en preparación)).

\item Análisis empírico de equidad por fototipo Fitzpatrick sobre
Fitzpatrick17k, comparando cinco arquitecturas y reportando la
brecha de rendimiento entre fototipos extremos.

\item Caracterización de la sensibilidad de DermLIP al tokenizador
y al diseño del \emph{prompt} en clasificación zero-shot, con
identificación del factor crítico (compatibilidad del tokenizador
con el espacio de \emph{embeddings} del modelo)~(Anexo \emph{Comparación LLM extendida} (en preparación)).

\item Replicación parcial del mecanismo de interpretabilidad
basado en Sparse Autoencoders sobre las representaciones de PanDerm
Large, validado mediante Concept Bottleneck Model sobre
SkinCon~\cite{skincon} (Anexo \emph{SAE y conceptos clínicos} (en preparación)).

\item Prototipo funcional DermApIxel como entorno integrado de
despliegue de los módulos derivados del trabajo
(Anexo \emph{Prototipo DermApIxel} (en preparación)).
\end{enumerate}

\section{Restricciones del estudio}
\label{sec:alcance}

Tres restricciones explícitas condicionan el diseño experimental y
la interpretación de los resultados. \textbf{Hardware:} todas las
ejecuciones se realizan sobre NVIDIA DGX Spark con chip Grace
Hopper GB10 y arquitectura aarch64, no coincidente con el hardware
x86\_64 dominante en la literatura. \textbf{Acceso a datos:} los
conjuntos internos del grupo de Monash empleados en el
preentrenamiento de PanDerm no son accesibles para terceros, lo que
limita el trabajo a los pesos liberados y a la evaluación sobre
datasets externos públicos. \textbf{Acceso a modelos:} los pesos
de DermFM-Zero no son públicos en la fecha de cierre, lo que impide
su evaluación empírica directa y restringe su tratamiento a la
discusión conceptual.

El trabajo no aborda, por decisión de alcance, los siguientes
elementos: preentrenamiento desde cero de un modelo fundacional
propio; validación clínica prospectiva sobre cohorte hospitalaria
real con protocolo aprobado por Comité Ético de Investigación
Clínica; integración operativa con el sistema PACS del HUSLL;
reentrenamiento multilingüe de la rama lingüística sobre datasets en
español de gran escala. Estos elementos se identifican como líneas
futuras en la apartado \emph{Líneas de investigación abiertas} (en preparación) del
capítulo \emph{Conclusiones} (en preparación).

\section{Estructura del documento}
\label{sec:estructura}

El documento se compone de siete capítulos en el cuerpo principal,
una bibliografía y diez anexos. El capítulo 2 sintetiza
el estado del arte en cuatro bloques. El
capítulo \emph{Materiales y métodos} (en preparación) describe materiales y métodos: los doce
datasets externos, la ontología jerárquica L1/L2/L3, el dataset
DermapixelAI~1.0 (resumido; pipeline detallado en el
Anexo \emph{Pipeline DermapixelAI} (en preparación)), los catorce modelos evaluados, la
infraestructura, los protocolos, las métricas y la trazabilidad. El
capítulo \emph{Resultados} (en preparación) reúne los resultados experimentales
organizados por bloque. El capítulo \emph{Discusión} (en preparación) interpreta
los resultados en clave de hallazgos, comparación con la literatura
e implicaciones operativas. El capítulo \emph{Limitaciones} (en preparación)
declara las limitaciones del estudio en seis dimensiones. El
capítulo \emph{Conclusiones} (en preparación) responde de forma compacta a la
pregunta de investigación y enumera las contribuciones
reproducibles.

Los anexos contienen, respectivamente, el mapa de tareas
experimentales y la estructura del repositorio
(Anexo \emph{Tareas y reproducibilidad} (en preparación)), las tablas detalladas de resultados por
dataset (Anexo \emph{Tablas extendidas} (en preparación)), el pipeline de construcción y
caracterización del dataset DermapixelAI~1.0
(Anexo \emph{Pipeline DermapixelAI} (en preparación)), la declaración sobre el uso de
herramientas de inteligencia artificial durante la elaboración del
documento (Anexo \emph{Declaración uso IA} (en preparación)), el documento de entrega formal
del dataset con licencia, datasheet y consideraciones de privacidad
(Anexo \emph{Entrega del dataset} (en preparación)), la documentación de los Sparse
Autoencoders y el diccionario de conceptos
(Anexo \emph{SAE y conceptos clínicos} (en preparación)), la comparación extendida con modelos
de lenguaje multimodales generalistas (Anexo \emph{Comparación LLM extendida} (en preparación)) y
el prototipo DermApIxel (Anexo \emph{Prototipo DermApIxel} (en preparación)).
